\documentclass{article}
\usepackage{algorithm, algorithmic, amsfonts, amsmath}

\begin{document}
\begin{algorithm}[!h]
\renewcommand{\algorithmicrequire}{\textbf{Input:}}
\renewcommand{\algorithmicensure}{\textbf{Output:}}
	\caption{PUL-SITS: PU learning for satellite image time series}
	\begin{algorithmic}[1]
	    \REQUIRE
	    Positive time series $P$; unlabeled set $U$; Huber threshold $\delta$;
	    consistency weight $\lambda$; epochs $E_{ae},E_{clf}$; batch size $B$
	    \ENSURE
	    Trained binary classifier $f_{\theta}$
	    \STATE \textbf{Stage 1: select reliable negatives}
	    \STATE Train a variational recurrent autoencoder $VAERNN$ (two-layer GRU)
	    \emph{on $P$ only}, minimising the Huber reconstruction loss % having seen positives only, it reconstructs them well and everything else poorly
	    \STATE $E(TS)\leftarrow HLoss\big(TS,\ \mathrm{Linear}(VAERNN(TS))\big)$ for each $TS\in U$
	    \STATE $\mu_{rec}\leftarrow \frac{1}{|U|}\sum_{TS\in U}E(TS)$;\quad
	    $U_{\mu}\leftarrow\{TS\in U\mid E(TS)>\mu_{rec}\}$ % a reconstruction error above the mean marks a large difference from the positive class
	    \STATE $RN\leftarrow$ a random subset of $U_{\mu}$ with $|RN|=|P|$;\quad
	    $\tilde{U}\leftarrow U\setminus U_{\mu}$ % the low-error remainder becomes the unlabeled pool for consistency regularization
	    \STATE \textbf{Stage 2: two-model semi-supervised training}
	    \STATE Initialise identical FCGRU classifiers $f_{\theta}$ (main) and $f_{\hat{\theta}}$ (auxiliary)
	    \STATE Build reconstructed copies
	    $\hat{P}=\{\mathrm{Linear}(VAERNN(TS))\}_{TS\in P}$ and $\hat{RN}$ likewise;
	    label $P,\hat{P}$ as $1$ and $RN,\hat{RN}$ as $0$ % the autoencoder serves a second purpose: producing perturbed versions
	    \FOR{$epoch=1$ \TO $E_{clf}$}
	    \FORALL{steps}
	    \STATE Train $f_{\hat{\theta}}$ by BCE on a batch of the \emph{reconstructed} labelled data
	    $\hat{B}\subset\hat{P}\cup\hat{RN}$
	    \STATE Sample $B_{l}\subset P\cup RN$ and $B_{u}\subset\tilde{U}$;
	    reconstruct $\hat{TS}_{j}$ for each $TS_{j}\in B_{u}$
	    \STATE Update $f_{\theta}$ on
	    $\mathcal{L}_{BCE}(B_{l};f_{\theta})
	    +\lambda\cdot\frac{1}{|B_{u}|}\sum_{j}\mathrm{KL}\!\left(f_{\hat{\theta}}(\hat{TS}_{j})\,\|\,f_{\theta}(TS_{j})\right)$ % the main model sees the original data and the auxiliary model the reconstruction; their predictions should agree
	    \ENDFOR
	    \ENDFOR
		\RETURN $f_{\theta}$
	\end{algorithmic}
\end{algorithm}
% Desloires, Ienco, Botrel, Remote Sensing 2022, "Positive Unlabelled Learning for
% Satellite Images' Time Series Analysis: An Application to Cereal and Forest Mapping".
% Task: in remote-sensing mapping only one land-cover class has field samples and every other pixel is unlabeled,
% which is the classic PU structure. The key design is that the autoencoder does two jobs:
% (1) it selects reliable negatives. Trained on P alone, it reconstructs positives well and other classes poorly, so
%     anything with a reconstruction error above the mean is a candidate negative;
% (2) it generates perturbations. The reconstructed version acts as another view of the same sample for the
%     consistency regularizer of the second stage.
% The two halves of the data serve different ends: the high-error half of U (U_mu) supplies RN, while the low-error
% half (U~) becomes the unlabeled pool for consistency regularization. The latter is likelier to contain true
% positives, which is exactly the part the consistency constraint is there to exploit. The two models divide the work:
% the auxiliary trains on reconstructions, the main model on originals, and a KL divergence holds them in agreement.
\end{document}
